\documentclass[11pt]{article}

\usepackage[T1]{fontenc}
\usepackage[utf8]{inputenc}
\usepackage{lmodern}
\usepackage[margin=0.75in]{geometry}
\usepackage{xcolor}
\usepackage{booktabs}
\usepackage{tabularx}
\usepackage{array}
\usepackage{listings}
\usepackage{hyperref}
\usepackage{titlesec}
\usepackage{enumitem}
\usepackage[most]{tcolorbox}
\usepackage{microtype}

\definecolor{accent}{HTML}{0B6E69}
\definecolor{accentdark}{HTML}{084C49}
\definecolor{ink}{HTML}{1E2933}
\definecolor{muted}{HTML}{5B6770}
\definecolor{soft}{HTML}{F2F7F6}
\definecolor{line}{HTML}{C9D8D5}
\definecolor{codebg}{HTML}{F7F9FA}

\hypersetup{
  colorlinks=true,
  linkcolor=accentdark,
  urlcolor=accentdark,
  citecolor=accentdark,
  pdftitle={RTX 5090 FP8/FP6/FP4 Reduction Tree Width Report},
  pdfauthor={tensor-cores-numerical-behavior}
}

\titleformat{\section}
  {\Large\bfseries\color{accentdark}}
  {\thesection}{0.55em}{}
\titleformat{\subsection}
  {\large\bfseries\color{accentdark}}
  {\thesubsection}{0.55em}{}
\titlespacing*{\section}{0pt}{1.3em}{0.55em}
\titlespacing*{\subsection}{0pt}{0.9em}{0.35em}

\setlist[itemize]{leftmargin=1.2em, itemsep=0.2em, topsep=0.25em}
\setlength{\parindent}{0pt}
\setlength{\parskip}{0.55em}
\renewcommand{\arraystretch}{1.22}

\lstdefinestyle{reportcode}{
  basicstyle=\ttfamily\footnotesize,
  backgroundcolor=\color{codebg},
  frame=single,
  rulecolor=\color{line},
  framerule=0.5pt,
  breaklines=true,
  columns=fullflexible,
  keepspaces=true,
  showstringspaces=false,
  tabsize=2
}
\lstset{style=reportcode}

\newcolumntype{Y}{>{\raggedright\arraybackslash}X}
\newcolumntype{R}{>{\raggedleft\arraybackslash}p{0.72in}}
\newcolumntype{C}{>{\centering\arraybackslash}p{0.72in}}

\newtcolorbox{summarybox}{
  enhanced,
  colback=soft,
  colframe=accent,
  boxrule=0.7pt,
  arc=2mm,
  left=8pt,
  right=8pt,
  top=8pt,
  bottom=8pt
}

\newtcolorbox{metricbox}[2]{
  enhanced,
  colback=white,
  colframe=line,
  boxrule=0.6pt,
  arc=2mm,
  left=7pt,
  right=7pt,
  top=6pt,
  bottom=6pt,
  title={#1},
  coltitle=accentdark,
  fonttitle=\bfseries,
  attach boxed title to top left={xshift=5pt,yshift=-2mm},
  boxed title style={colback=soft,colframe=line,boxrule=0.4pt,arc=1mm},
  #2
}

\begin{document}
\pagenumbering{gobble}

\begin{titlepage}
\vspace*{0.4in}
{\Huge\bfseries\color{accentdark} RTX 5090 FP8/FP6/FP4\\ Reduction Tree Width Report\par}
\vspace{0.2in}
{\Large Direct PTX tensor-core boundary measurements\par}
\vspace{0.25in}
{\color{muted}Date: 2026-05-09\par}
\vfill
\begin{summarybox}
\textbf{Headline finding.}
All tested FP8, FP6, and FP4 direct-PTX QMMA variants expose a
\textbf{21-bit effective addition/reduction boundary} for the
\texttt{m16n8k32} probe. The binary32 reference for the same probe is
\textbf{24 bits}, so the measured low-precision tensor-core path is
\textbf{3 effective bits narrower} in this test.
\end{summarybox}
\vspace{0.2in}
\begin{tabularx}{\linewidth}{>{\bfseries}l Y}
\toprule
GPU & NVIDIA GeForce RTX 5090, compute capability 12.0 \\
CUDA & Toolkit 13.2, \texttt{nvcc} V13.2.78 \\
Build target & \texttt{compute\_120a} / \texttt{sm\_120a} \\
Source & \texttt{src/tc\_test\_numerics-5090-lowp-reduction-width.cu} \\
\bottomrule
\end{tabularx}
\end{titlepage}
\clearpage
\pagenumbering{arabic}

\section{Scope}

The tests execute inline PTX \texttt{mma.sync.aligned} instructions directly
from a single warp. They do not call cuBLASLt or any GEMM library. The reported
quantity is an externally observable effective addition/reduction width for one
\texttt{m16n8k32} instruction. It should not be read as a complete physical map
of the internal reduction tree.

FP8 uses regular unscaled FP8 MMA instructions. FP6 and FP4 use the PTX
\texttt{kind::f8f6f4} form. Compact block-scaled MXFP4/MXFP6 instructions are
outside this report.

\section{Methodology}

Each measurement runs one direct tensor-core MMA:

\[
  D = C + \sum_{i=0}^{31} A_i B_i
\]

All A and B elements are encoded as exactly \texttt{1.0}. The exact dot-product
contribution is therefore \(K=32\). The C accumulator is a binary32 value set
to \(2^{shift}\). The probe scans shifts from 16 to 40 and records the largest
shift where the output is still greater than the base accumulator.

\begin{summarybox}
\[
  p = max\_changed\_shift - \log_2(K) + 1
\]
\[
  K = 32,\quad \log_2(K)=5,\quad max\_changed\_shift=25,\quad p=21
\]
\end{summarybox}

A full binary32 significand has 24 bits. For the same \(K=32\) boundary probe,
the binary32 reference max shift is:

\[
  \log_2(32) + 24 - 1 = 28
\]

\section{Key Code Snippets}

The constants define the K dimension, scan interval, and binary32 reference.

\begin{lstlisting}[language=C++]
static const int kInstructionK = 32;
static const int kLog2InstructionK = 5;
static const int kMinShift = 16;
static const int kMaxShift = 40;
static const int kExpectedBinary32Bits = std::numeric_limits<float>::digits;
static const int kExpectedBinary32MaxShift =
    kLog2InstructionK + kExpectedBinary32Bits - 1;
\end{lstlisting}

FP4 E2M1 values are shifted into the central four bits required by the
\texttt{kind::f8f6f4} PTX operand format.

\begin{lstlisting}[language=C++]
static std::uint8_t encode_fp4_e2m1_padded(float value) {
  return static_cast<std::uint8_t>(__nv_fp4_e2m1(value).__x << 2);
}
\end{lstlisting}

The MMA path is direct inline PTX. The output operands are binary32 registers.

\begin{lstlisting}[language=C++]
#define TCNB_MMA_ASM(instruction)                                               \
  asm volatile(instruction " {%0, %1, %2, %3}, "                                \
                           "{%4, %5, %6, %7}, {%8, %9}, "                      \
                           "{%10, %11, %12, %13};\n"                           \
               : "=f"(d0), "=f"(d1), "=f"(d2), "=f"(d3)                       \
               : "r"(a_pack), "r"(a_pack), "r"(a_pack), "r"(a_pack),         \
                 "r"(b_pack), "r"(b_pack), "f"(base), "f"(base),             \
                 "f"(base), "f"(base))
\end{lstlisting}

The scan records the last visible contribution and converts it to bits.

\begin{lstlisting}[language=C++]
for (int shift = kMinShift; shift <= kMaxShift; ++shift) {
  float base = std::ldexp(1.0f, shift);
  float sample = run_once<variant>(a_pack, b_pack, shift, device_out);

  if (sample > base) {
    result.max_changed_shift = shift;
    result.sample_at_max = sample;
  } else if (result.max_changed_shift >= 0 &&
             result.sample_after_max == 0.0f) {
    result.sample_after_max = sample;
  }
}

result.observed_bits =
    result.max_changed_shift - kLog2InstructionK + 1;
\end{lstlisting}

The build targets the accelerated SM120a feature set explicitly.

\begin{lstlisting}[language=make]
test-5090-fp8-reduction-width: $(SRC_5090_LOWP_REDUCTION_WIDTH)
	$(NVCC) $(NVCC_INCLUDES) -DTCNB_REDUCTION_FP8 -o $@ \
	  -arch=$(GPU_COMPUTE_5090_ACCEL) -code=$(GPU_SM_5090_ACCEL) \
	  $(NVCC_STD) $<
\end{lstlisting}

\section{Instructions Tested}

{\footnotesize
\begin{tabularx}{\linewidth}{l l Y}
\toprule
Family & Variant & PTX instruction \\
\midrule
FP8 & E4M3 x E4M3 & \texttt{mma.sync.aligned.m16n8k32.row.col.f32.e4m3.e4m3.f32} \\
FP8 & E5M2 x E5M2 & \texttt{mma.sync.aligned.m16n8k32.row.col.f32.e5m2.e5m2.f32} \\
FP8 & E4M3 x E5M2 & \texttt{mma.sync.aligned.m16n8k32.row.col.f32.e4m3.e5m2.f32} \\
FP8 & E5M2 x E4M3 & \texttt{mma.sync.aligned.m16n8k32.row.col.f32.e5m2.e4m3.f32} \\
FP6 & E2M3 x E2M3 & \texttt{mma.sync.aligned.m16n8k32.row.col.kind::f8f6f4.f32.e2m3.e2m3.f32} \\
FP6 & E3M2 x E3M2 & \texttt{mma.sync.aligned.m16n8k32.row.col.kind::f8f6f4.f32.e3m2.e3m2.f32} \\
FP4 & E2M1 x E2M1 & \texttt{mma.sync.aligned.m16n8k32.row.col.kind::f8f6f4.f32.e2m1.e2m1.f32} \\
\bottomrule
\end{tabularx}
}

\section{Results}

\begin{metricbox}{Core numbers}{}
\begin{tabularx}{\linewidth}{>{\bfseries}l Y}
\toprule
Largest shift where \(D>C\) & 25 \\
Observed effective width & 21 bits \\
Binary32 reference max shift & 28 \\
Binary32 reference width & 24 bits \\
Gap from binary32 reference & -3 bits \\
\bottomrule
\end{tabularx}
\end{metricbox}

\begin{tabularx}{\linewidth}{l l R R R R}
\toprule
Family & Variant & Max shift & Sample at max & First unchanged & Bits \\
\midrule
FP8 & E4M3 x E4M3 & 25 & 33,554,464 & 67,108,864 & 21 \\
FP8 & E5M2 x E5M2 & 25 & 33,554,464 & 67,108,864 & 21 \\
FP8 & E4M3 x E5M2 & 25 & 33,554,464 & 67,108,864 & 21 \\
FP8 & E5M2 x E4M3 & 25 & 33,554,464 & 67,108,864 & 21 \\
FP6 & E2M3 x E2M3 & 25 & 33,554,464 & 67,108,864 & 21 \\
FP6 & E3M2 x E3M2 & 25 & 33,554,464 & 67,108,864 & 21 \\
FP4 & E2M1 x E2M1 & 25 & 33,554,464 & 67,108,864 & 21 \\
\bottomrule
\end{tabularx}

At shift 25, the base is \(2^{25}=33{,}554{,}432\), and the observed sample is
33,554,464. The \(+32\) contribution is still visible. At shift 26, the first
unchanged sample is 67,108,864, equal to \(2^{26}\). The contribution is no
longer visible.

\section{Facts and Evidence}

\begin{tabularx}{\linewidth}{>{\bfseries}p{1.75in} Y}
\toprule
Fact & Evidence \\
\midrule
Direct PTX path & The source emits inline \texttt{mma.sync.aligned}; no cuBLASLt dependency is present. \\
Single-instruction boundary & Each sample launches one warp and one \texttt{m16n8k32} MMA for the selected variant. \\
Exact input contribution & All A and B operands encode \texttt{1.0}; the K dimension contributes exactly 32. \\
Uniform observed width & All FP8, FP6, and FP4 result files report max shift 25 and width 21. \\
Binary32 comparison & Binary32 would produce max shift 28 for the same K and boundary formula. \\
\bottomrule
\end{tabularx}

\subsection{SASS Confirmation}

{\footnotesize
\begin{tabularx}{\linewidth}{l Y}
\toprule
Target & Confirmed SASS mnemonics \\
\midrule
\texttt{test-5090-fp8-reduction-width} & \texttt{QMMA.16832.F32.E4M3.E4M3}, \texttt{QMMA.16832.F32.E5M2.E5M2}, \texttt{QMMA.16832.F32.E4M3.E5M2}, \texttt{QMMA.16832.F32.E5M2.E4M3} \\
\texttt{test-5090-fp6-reduction-width} & \texttt{QMMA.16832.F32.E2M3.E2M3}, \texttt{QMMA.16832.F32.E3M2.E3M2} \\
\texttt{test-5090-fp4-reduction-width} & \texttt{QMMA.16832.F32.E2M1.E2M1} \\
\bottomrule
\end{tabularx}
}

\section{Reproduction}

\begin{lstlisting}[language=bash]
make -B test-5090-fp8-reduction-width \
        test-5090-fp6-reduction-width \
        test-5090-fp4-reduction-width

python3 run_tests.py -o 5090 \
        5090-fp8-reduction-width \
        5090-fp6-reduction-width \
        5090-fp4-reduction-width

cuobjdump --dump-sass test-5090-fp8-reduction-width
cuobjdump --dump-sass test-5090-fp6-reduction-width
cuobjdump --dump-sass test-5090-fp4-reduction-width
\end{lstlisting}

\section{Interpretation}

The root cause of the 24-bit binary32 expectation failure is the low-precision
QMMA accumulation/reduction path itself. The inputs are exact, the dot-product
contribution is exact, and SASS confirms direct tensor-core QMMA instructions.
The measured path simply stops preserving the \(+32\) contribution three shifts
earlier than a full binary32 significand path would.

The result is consistent across all tested FP8, FP6, and FP4 input type
combinations: the observable effective addition/reduction boundary is 21 bits.

\section{Limitations}

\begin{itemize}
\item This is a single-instruction \texttt{m16n8k32} boundary probe, not a full
GEMM accuracy characterization.
\item The measurement reports observable effective width, not the internal
physical topology of the hardware reduction tree.
\item FP6 and FP4 are tested through the PTX \texttt{kind::f8f6f4} form.
Compact block-scaled MXFP4/MXFP6 instructions are not covered.
\item The input pattern is all ones. Other value distributions can exercise
different rounding and normalization behavior.
\end{itemize}

\section{Summary}

\begin{summarybox}
RTX 5090 direct-PTX FP8, FP6, and FP4 tensor-core MMA instructions tested here
all measure at 21 effective reduction/addition bits for the \texttt{m16n8k32}
boundary probe. The binary32 reference is 24 bits. The observed low-precision
QMMA path is therefore consistently 3 effective bits narrower than binary32 for
this measurement.
\end{summarybox}

\end{document}
